% ======================================================================
% Ontologie der Kontinua -- Core 1.3
% Cross-K-Modul: crossk_k7_k8.tex
% Cross-Level-Relationen zwischen K7 und K8
% ======================================================================

\subsubsection{K7–K8-Querverbindung}
\label{sec:crossk-k7-k8}

Der Übergang $K_7 \rightarrow K_8$ markiert die Entstehung des
zivilisatorisch-technologischen Kontinuums aus dem sozialen Kontinuum.
Während $K_7$ Normen, Rollen, Kommunikationsflüsse und institutionelle
Zyklen bereitstellt, führt $K_8$ ein:
\begin{itemize}
  \item technologische und infrastrukturelle Achsen,
  \item großskalige symbolische und informationelle Systeme,
  \item Energie- und Logistiknetze,
  \item ökonomische Reproduktionsmechanismen,
  \item Reparatur-, Wartungs- und höherstufige Kooperationszyklen,
  \item globale Kontinuitätsbedingungen für populationsweite Systeme.
\end{itemize}

$K_8$ ist das erste Kontinuum, das stabile Strukturen tragen kann, die
die kognitive und soziale Reichweite einzelner Akteure übersteigen.

% ----------------------------------------------------------------------
\subsubsection{1. Stufen und ihre Rollen}

\paragraph{K7 (soziales Kontinuum).}
$K_7$ liefert:
\begin{itemize}
  \item soziale Normen und Rollen ($A_{\mathrm{soc}}$),
  \item Kommunikationsflüsse $J_{\mathrm{comm}}$,
  \item institutionelle Zyklen $C_{\mathrm{inst}}$,
  \item kollektive Potentiale $P_{\mathrm{soc}}$ (Vertrauen, Legitimität),
  \item Kohäsions- und Durchsetzungsschwellen $\Theta_{\mathrm{soc}}$.
\end{itemize}

Allerdings fehlen in $K_7$:
\begin{itemize}
  \item großskalige Infrastrukturen,
  \item technologische Produktionszyklen,
  \item Energie-Dichteschwellen,
  \item materielle und informationelle Logistik,
  \item stabile multigenerationale Reproduktionsmechanismen.
\end{itemize}

\paragraph{K8 (zivilisatorisch-technologisches Kontinuum).}
$K_8$ führt ein:
\begin{itemize}
  \item Achsen für Infrastruktur, Technik und symbolische Systeme ($A_{\mathrm{tech}}$),
  \item Potentiale für Energiedichte, Reparaturkapazität, Logistikdurchsatz,
  \item Flüsse $J^{(8)} = (J_{\mathrm{energy}}, J_{\mathrm{material}},
        J_{\mathrm{logistics}}, J_{\mathrm{info}}, J_{\mathrm{population}})$,
  \item Infrastrukturzyklen $C_{\mathrm{sys}}$ (Produktion → Transport → Nutzung → Wartung),
  \item Schwellen $\Theta_8 = (\Theta_E, \Theta_L, \Theta_I, \Theta_M, \Theta_R)$,
  \item Maß $k_8$ für zivilisatorische Kontinuität.
\end{itemize}

Damit ist $K_8$ das erste Kontinuum mit großskaliger energetischer und
informationeller Autonomie.

% ----------------------------------------------------------------------
\subsubsection{2. Gemeinsame und vererbte Achsen sowie Schwellen}

\paragraph{Geteilte symbolische Strukturen.}
Symbolische Strukturen und Kommunikationssysteme von $K_7$ werden zum Substrat für:
\[
A_{\mathrm{symbol}}(K_7) \rightarrow A_{\mathrm{tech}}(K_8),
\]
was technische Kodifikation ermöglicht (Schrift, Mathematik, Engineering,
Standards).

\paragraph{Neue Achsen in \texorpdfstring{$K_8$}{K_8}.}
Dazu zählen:
\begin{itemize}
  \item $A_{\mathrm{infra}}$ — Straßen, Stromnetze, Wassersysteme,
  \item $A_{\mathrm{energy}}$ — Energiedichte und Umwandlungssysteme,
  \item $A_{\mathrm{econ}}$ — ökonomische Reproduktionsmechanismen,
  \item $A_{\mathrm{info}}$ — Informationsspeicherung und -übertragung,
  \item $A_{\mathrm{repair}}$ — Wartung und Regeneration.
\end{itemize}

\paragraph{Schwellen.}
Vererbt bleiben Kohäsionsschwellen aus $K_7$, sind aber unzureichend.
Neue Schwellen entstehen:
\begin{itemize}
  \item $\Theta_E$ — minimale Energiedichte für systemische Reproduktion,
  \item $\Theta_L$ — Logistikdurchsatzschwellen,
  \item $\Theta_I$ — informationelle Kohärenzschwellen,
  \item $\Theta_M$ — Reparaturschwelle,
  \item $\Theta_R$ — demografische Nachhaltigkeitsschwelle.
\end{itemize}

Falls eine dieser Schwellen versagt, kollabiert $\Omega(K_8)$.

% ----------------------------------------------------------------------
\subsubsection{3. Cross-Level-Flüsse, Zyklen und Spannungen}

\paragraph{Flüsse.}
Flüsse in $K_8$ verallgemeinern und skalieren soziale Flüsse:
\[
J^{(8)} = (J_{\mathrm{energy}}, J_{\mathrm{material}},
          J_{\mathrm{logistics}}, J_{\mathrm{info}},
          J_{\mathrm{population}}).
\]

\paragraph{Zyklen.}
Ziv.-technologische Stabilität erfordert den Infrastrukturzyklus $C_{\mathrm{sys}}$:
\[
\text{Produktion} \rightarrow \text{Transport} \rightarrow
\text{Nutzung} \rightarrow \text{Wartung} \rightarrow \text{Produktion}.
\]
Stabilitätsbedingung:
\[
\oint dR \approx 0.
\]

\paragraph{Spannung.}
Cross-Level-Spannung:
\[
T_{7,8} =
f(\Theta_E - P_{\mathrm{energy}},\
  \Theta_L - J_{\mathrm{logistics}},\
  \Theta_I - P_{\mathrm{info}},\
  \Theta_M - P_{\mathrm{repair}},\
  \Theta_R - P_{\mathrm{population}}).
\]

Übermäßige Spannung zerstört infrastrukturelle Kontinuität.

% ----------------------------------------------------------------------
\subsubsection{4. Geburts- und Todesbedingungen über die Stufen}

\paragraph{Geburt von \texorpdfstring{$K_8$}{K_8}.}
$K_8$ entsteht, wenn:
\[
T_7 > \Theta_{7,\mathrm{dim}}
\quad\text{und}\quad
A_{\mathrm{tech}}, A_{\mathrm{infra}} \in M_8 \setminus A(K_7).
\]

Die Domäne erweitert sich:
\[
\Omega(K_8) = \Omega(K_7) \cup \Delta\Omega_{\mathrm{civ}}.
\]

\paragraph{Todesausbreitung.}
Falls soziale Schwellen versagen ($\Theta_{\mathrm{soc}}$), verlieren
zivilisatorische Systeme Kohärenz und fallen auf $K_7$ zurück.

Falls eine der fünf $K_8$-Schwellen versagt:
\[
\Theta_E,\Theta_L,\Theta_I,\Theta_M,\Theta_R,
\]
dann gilt $\Omega(K_8)=\varnothing$.

$K_7$ kann stabil bleiben, sofern seine eigenen Schwellen nicht verletzt sind.

% ----------------------------------------------------------------------
\subsubsection{5. Konzeptionelle Beispiele}

\paragraph{Eskalation der Energiedichte.}
Wenn die Energiedichte $\Theta_E$ überschreitet, werden technische und
ökonomische Strukturen selbsttragend.

\paragraph{Informationskodifizierung.}
Soziale Symbole entwickeln sich zu technischen Systemen (Schrift → Mathematik →
Wissenschaft), die neue Achsen in $A_{\mathrm{tech}}$ ermöglichen.

\paragraph{Infrastrukturversagen.}
Fällt die Logistikkapazität unter $\Theta_L$ oder Reparaturraten unter
$\Theta_M$, brechen zivilisatorische Zyklen und das System bricht zusammen:
\[
C_{\mathrm{sys}} \nrightarrow \text{geschlossene Schleife}.
\]

Damit lässt sich der K7→K8-Übergang in komprimierter Form als Auftreten
technisch-ziviler Struktur aus sozialen Normen, Kommunikationsflüssen und
institutioneller Stabilität zusammenfassen.
